\chapter{The Ordonian Calculus and the IMS: Stewardship, Agency, and the Economy of Information in The Systems Experiment}

\begin{center}
    Author: Llamangus. Editor: Harlow. Peer reviewed. % clearly indicate non-peer reviewed articles!
\end{center}

\begin{flushright}
    IMSC: j(2)-ls-1

    DOI: \doilink{10.5281/zenodo.23702122}
\end{flushright}

% ============================================================================
\section*{Abstract}
Agency and Stewardship theories can be used to describe how we understand the relationship between delegates and delegators, both in the corporate world, in volunteer spaces, and in The Systems Event. We can understand how the lack of a clear framework, compounding issues from the organization of the event, and player siloing worked against the desired goals of the event. This essay draws on first-hand experiences and accounts from related essays, tracing how two competing teams diverged and were impacted during the events Exploration and Exploitation time frame, largely from the Ordonian perspective. “The Ordonian Calculus” is posed as a phenomenological analysis of social and systemic risks and anticipations, and their effects on team cohesion.

The essay also proposes an Agency Theory of Minecraft, arguing that the default singleplayer design leaves multiplayer environments particularly exposed to agentic, exploitative play absent of deliberate design intervention, and offers guidance for future IMS events.

\section{Introduction}
We were down to `the core four', the burnt-out remnants of the presidium by the third week of The Systems Experiment (TSE): Pheelix(`Zandercamper'), Blaze (`blAze\_floW'), Tiff (`Emo\_Cupcakes'), and myself – Llamangus (`ComradeLlamangus'). The Presidium was Pheelix's design, intended to be a temporary structure for the earliest days of the two-month long exploitation period.\footnote{Or \textit{Exploitation phase} in the first Systems Event. The execution and implementation of systems and the realization of their value within the existing structures and bodies of knowledge.} It was less than halfway there when even the core four began to log in less and less.

Along with designing the presidium, Pheelix created the Discord and test server, and was the proactive leader that Ordonia needed in the exploration phase.\footnote{Or \textit{Exploration phase} in the first Systems Event. The generation of novel assumptions and ideas to be proven useful} Blaze created a detailed, towering city design with towers and aqueducts. Tiff, perhaps our most active player, created most of the farms we could theoretically need for our systems.

Ultimately, we discovered a lot of our systems were overbuilt for the conditions we found ourselves in, as we were prepared for a team of ten to twenty. The Ordonian government never reached its potential, the city planning was slow and inflexible, and there wasn't much of an economy to build around.

I addressed the presidium in the \#leadership-private section of our Discord as there were things that needed to be said. The event had not gone as planned, and there was no reason to expect that conditions of the event would change to meet the needs of our designs. We each had to pick up increasing amounts of responsibilities as most of our team had either been no-shows on the server or had dropped out. None of the core four wanted to be seen as the group's failure point, but each of us at one point or another in the event had taken on large tasks. In Rev's work titled `Who Decides? Leadership, Legitimacy, and Load in competing teams', he pointed out during interviews that members' perception of the shared workload differed fairly dramatically, depending on what systems they were in charge of.~\citep{Rev(2026)}

It seemed like the only option was to reorganize, learn to communicate more effectively, and reassess our goals in this event. By this point even the organizers knew this event would not produce the desired outcomes. However, even `bad' Minecraft, like pizza, is still `good Minecraft'.

The next day, Turtle announced that the event would end and be turned into an SMP with the opening of the End and the dragon fight the following week, but even the SMP would shut down for good less than a week later as the Institute shifted its focus for August. It was a relief that it was allowed to end on a good note.

In light of the failures of the Systems Experiment, guidelines need to be set in place by the IMS in order to clear the ambiguity of players, the IMS needs to understand what resources they have as far as players and players time, and Agency and Stewardship\footnote{a model of behavior wherein an individual is not simply motivated by individualistic utility, but is instead a steward whose motives are aligned with those of the Principal and/or collective, even when those interests may diverge} should be adopted as a framework for understanding the social underpinnings of servers, especially in regard to leadership due to the similarities between server event planning and a business.

\section{The setup}
According to early promotional materials for the event, it is described as a team-based event where players are individually tasked with designing a system: “A repeatable structure of Rules, Incentives, or infrastructure.” These systems would need to fit one or more of the categories listed: Economic, Government, Society, Infrastructure, Incentive, Industry, or Defense. Players had roughly three weeks through the month of June to submit a systems draft of their idea for the event, with late comers given a grace period.

In Rev's article on the Systems Experiment, he describes the event in two phases: “Exploration” and “Exploitation”.~\citep{Rev(2026)} The promotional material calls this Exploration phase the `preseason'. This work will continue using Rev's reference.

During the event's Exploration phase, players were free to ask questions about how the event would be run. These questions were primarily answered by Turtle, the creator of TSE. One of the questions from `Dinosawrr' (dino.sawrr), as paraphrased by Turtle, asked, “can outside, non-systems players join as citizens?”~\citep{Turtle(2026a)} Turtle's response gives us some insight into his thought process as an event organizer, stating, “no, people cannot join without following the normal requirements and expectations”.~\citep{Turtle(2026a)} He goes on to explain:

\begin{quotation}
    While it sounds fine in theory, this will introduce a lot of members who want to focus specifically on SMP style play. They are likely to get annoyed with systems over time and bypass them. In addition, we want systems to use each other (i.e., industry system player sells goods to infrastructure player with a currency from an economic system player), which might be diluted by non systems players.

    I believe it is entirely possible to design the systems you want without non systems players, and I think we should think about it more like a cross between Hermitcraft and Ish civ than just a civ. You people will have more time than the average team, you will be closer, and you will design systems that are more intricate. Besides, Minecraft is, as we have seen, a very easy game for single/low player counts.~\citep{Turtle(2026a)}
\end{quotation}

This shows that Turtle had some foresight into the problems that The Systems Event could, and in fact, did run into. From a first glance, it's clear that Turtle and the organizers at the IMS wanted this event to produce a certain kind of outcome. The Systems `not an experiment' Experiment was created to develop and showcase emergent systems. These systems would be developed competitively by two teams, Alpha and Beta (Beta would later be renamed to Ordonia).

Two examples were loosely cited as inspiration for this event. Civilization Event servers, with a clear factional focus, large and equally distributed player base, and sessioned play. The IMS, should they choose something like this, must consider these requirements in advance of team formation. These kinds of events are well known to many in the IMS, with members who have watched or participated in these kinds of events.

The second example was Hermitcraft. It may not be realistic for the IMS to host events that resemble Hermitcraft per se, but there could be an organizational philosophy based on that model, except for a structure suitable for our experiments and events. This assumes that stewardship is valued over agentic gameplay,\footnote{Gameplay that optimizes for maximal utility of its systems as permitted by the game itself or playing as if crossing off a checklist. “Given the opportunity, players will optimize the fun out of a game.” A player engaged in Agentic Gameplay defaults to Exploitation of known-optimal strategies towards a model of complete-ness, often over the game's fuller possibility space. This mirrors the Agents pursuit of individual utility over the Principal or collective's broader interests.} meaning that robust systems, narratives, and shared lore are created with each member potentially being a force-multiplier for each other. In a server, this might mean looser teams, districts, and less focus on `social victories'. This doesn't mean an end to team-based competitions entirely, but it could be of a smaller scope inside another event, department, or SMP.

The setup was meant to be in a fairly `archaic' or `vanilla' environment, without much attention to balanced gameplay with the exception of the banning of End Crystals and Elytras.~\citep{Turtle(2026b),Turtle(2026c)} There were few mods on the server, but we did use Simple Voice Chat and a `Wiretap' mod that provided a speaker and receiver. There was also a mod for copper rails; however, it did not have a large effect on gameplay, and its presence went largely unnoticed, only being used by Team Alpha between their settlement and their ice boat path to the shared neutral mushroom island.

Turtle had found a seed that at a first glance seemed to be perfect for this event. Team Alpha was given the seed a day earlier to pick between the northern or southern continent – an advantage that the organizers gave as a concession due to Alpha's difficulty recruiting and keeping active members up to that point. Alpha picked the northern continent, which had the only Pale Garden biome and the only ancient city that was found. Ordonia was given the southern continent, which had the largest Badlands and Mangrove Swamp Biomes. Ordonia's location would prove useful later for the Nether, being right next to a Fortress and being the closest to the only Bastion we would get in this event.

\section{Good intentions, flawed design}
We had a month to draft our systems and work with our teams to refine and plan their implementation. This `exploration phase' was quite generous on time, with the hopes that enough people would have drafted and refined the systems they would use in the exploitation phase. A `near-vanilla' experience, with few guidelines or restrictions beyond the siloing\footnote{Sorting of players into teams, often for competitive or antagonistic purposes. Encouraging agentic playstyles and rationale for the benefit of teams. (Factions, Nations, Clans, etc.)} was intended to allow more emergent systems to be developed.

Siloing the player base into teams is a common approach in competitive PvP events; however, those team are usually balanced in number, large enough to support the logistical needs of the team, players understand the focus of the event before they go in, and the victory conditions are often congruent with this scheme. This mismatch in focus on top of the `vanilla' or `archaic' framing of the event produced a bifurcation of focus among a small and uneven player base.

We can see how choices in the design and execution of The Server Experiment cumulatively might have led to the ending previously described. In “The Consequences of Poor Design” by Harlow, one of the fundamental issues posed was the lack of clear communication which led to uncertainty and burnout among sections of the player base.~\citep{Harlow(2026)} The problem of poor communication is not the sole or original critical flaw of the event, but it had several down-stream effects. We should ask: why was communication here poor?

The Systems Event was meant to develop systems that are the most broadly applicable or replicable inside of many types of servers; when we think of `vanilla', `archaic', or `anarchic' types of environments, where little gets in the way of player agency, we are mostly adapting to the structural systems given to us by Mojang. The event also created structural particularities that detracted from its focus. Given its siloed team-based nature and the a priori creation of levels of secrecy, the experiment became less about the face-value goals and more of an impromptu analysis of leadership and player sociology. By implying formal barriers, TSE organization made a choice that produced civ-like narratives over Hermitcraft-like emergent collaborative systems.

Hermitcraft is a model for a server of systems-creators in an otherwise mostly vanilla environment (data and resource packs aside). Each new member has been vetted after being known by their members for a while. A series might last between one to two years. An emergent system might take anywhere from a weekend to several months. The key difference is that hermits are not usually siloed. Each hermit is in some respects a force multiplier for each other.

\section{Through a lens of Agency and Stewardship}
In the world of business and finance, there is Agent theory and Stewardship theory. Agency, the more dominant theory, defines formal collaborative relationships between Principals\footnote{In business, the Principal is the shareholders who appoint executives as managers of a business/organization. In TSE, both the organizing body and the teams act as Principals over the players.} and Agents.\footnote{a utility-maximizing subordinate in a business relationship} The Principals delegate to Agents. It is in the best interests of Principals to create control schemes to minimize their risk of mismatched interests between parties. This poses a prisoner's dilemma between Principal and Agent, where cooperation is put up against individualistic utility maximization. This is in contrast to Stewardship models, which attempts to refute the pessimistic, control-oriented Agency framework. Instead, `agent' gets re-framed as `manager'.\footnote{In business, the individual(s) entrusted to execute the Principals objectives. May behave either as a self-serving Agent, or collective serving steward, depending on situational and psychological factors. In Minecraft, a player is a manager of oneself and the environment.} Psychologists and sociologists criticize Agency Theory\footnote{In business, control schemes must be used to reduce the risks posed by conflicting interests between Principals and Agents.} as having an incomplete understanding of organizational life and motivations. So, the question is: What causes interests to be aligned?~\citep{Davis(1997)}

The corporate-legal environment, of course, is very different from the Institute of Minecraft Studies or the Systems Experiment. Unlike a shareholder-executive relationship, no capital or legal obligations bind Principals and Agents. What I am offering by borrowing this frame is a method of analysis of server or event life. What Stewardship and Agency describe here is the uncertainty of the motives of both delegates and delegators, and how that uncertainty shapes anticipation and control-seeking behavior. This uncertainty exists whether the currency is dividends or the credibility of systems. Furthermore, you could extend the corporate examples to something closer to our world. Volunteer or open-source governance more broadly shows Agency and Stewardship theories into non-contractual settings; the IMS already does this through implementing a shared mission and democratic modes of leadership in departments. Understanding these relationships within servers and the IMS broadly could help us understand what is needed to shape and reinforce the intrinsic motivation that we rely on.

This, of course, brings us back to the IMS and the event. The IMS is a decentralized organization of volunteers and collaborators. Turtle and the IMS are the originating Principal from which a stewardship-oriented culture has naturally arrived. This Principal relationship also comes as a form of what \citet{Rev(2026)} calls “homophily”.

While Turtle's response to Dinosawrr's question was eventually overturned due to a lack of active players near the July 1st launch of the server (Exploitation Phase), we can establish the cultural practice of Stewardship are, stated or not, within the focus of this event, particularly with its relationship to the IMS. Seetloveletter, a member of the IMS Discord server who has extensive experience with Civ-style events and SMPs, stated as such that the typical IMS member would not be representative of the typical Civ event roster.~\citep{seetite(2026)}

Despite Turtle's criticism of `Content SMPs' or `Scripted SMPs' – noting that the events and player archetypes are made to appeal to a consumer outside of said SMPs – it is still the case with events that the agents or managers could be described as serving two principals, and a calculus of sorts had arrived to bridge the gaps where there were mismatched goals and understandings.~\citep{Turtle(2026),Turtle(nd)}

\section{Formative consequences}
\begin{verse}
    Scorvon: “Ordaina [sic] had a whole city built, banks schools etc, one week later they were bombed and everything was destoryed [sic];” “you cant win if ur systems are gone”.~\citep{scorevon(2026a),scorevon(2026b)}

    blAze\_floW: “And your victory will be empty if we never finish our systems in the first place”~\citep{blazeflow(2026)}

    Scorvon: “[a] win's a win”~\citep{scorevon(2026c)}
\end{verse}

A lack of proper communication and foresight given towards the principles of this event caused a conflict of interest in how participants saw the purpose of the event. For some, it was meant to be a sportsmanly development of systems, or a casual laboratory under the eyes of the shared Principal of the IMS. For others, the most obvious approach was more of a binary.

This split likely formed because of the differences in how the teams form through the course of the exploration and early exploitation phases. Ordonia had the full extent of time and the active support of early members of the presidium, particularly with Pheelix and Blaze. We had two large systems meant to encompass Ordonian life well before the server actually opened, meant to be the all-stars of the competition. We believed we were well prepared and advantaged during a time when Alpha was seen, as Turtle put it, `The Resistance'.

This framing seemed to hold even as membership requirements loosened to drive recruitment. I invited a friend of mine to the event; through which either by neglect or the organizers' providence, was assigned to Team Alpha. Ordonia had a spy in Alpha's ranks!

Alpha, on the other hand, rose through likely because their conditions favored a different approach. Their leadership was more centralized under a figurehead like Scorvon than later Phantom\_Boi, but their society was considerably less centralized. There were no democratic presidium or secretaries appointed via assembly as Pheelix had in mind for Ordonia.

Where Ordonia was waiting for a city design suitable enough, orderly enough, aesthetically harmonic enough, or complete enough in its design – Alpha's decentralized approach benefited their city. They had detailed builds, bases, social spaces like a tea party garden and a colosseum. Ordonia had a collection of cabin-like builds among farms; an industrial district posing as a town, supporting 4 players hoping that somehow new life would be put into the server by some manner, while we deliberated systems that were dreamt up before the server opened.

Alpha's approach came with some drawbacks though. There were complaints about the casual and careless use of wind charges, the death of all but one armadillo in an XP farm that was built out in the open. To us, there seemed to be a lack of cohesion among members of Alpha. A couple members confided that they were considering building away from their town. There was a similar sentiment in Ordonia before any headway was made on the city.

While Ordonia was building up its industry, Alpha crossed off a lot of the usual accomplishments before finishing the colosseum. They exhausted the trial chambers on their side and asked us for the right to use trial chambers on our side. They completed their colosseum and practiced PvP. But eventually their team would be largely inactive, except for the day of the Nether expansion and the Ender Dragon fight.

\section{Formation of a calculus}
There were signs at the server launch that The Systems Experiment had started off wrong. There wasn't much of an official launch. Some Ordonians weren't even aware when it launched. We were set into spectator mode to explore a map in which unbeknownst to us had a smaller world border than we had planned for, until we were able to convince Turtle to open the map up to its full intended extent. Later, most of the Ordonians that would be active were set to survival mode on the correct side of the map, but Team Alpha would continue to face issues with being in spectator or adventure mode until the next day. Some even spawned or respawned suffocating underground.

\subsection{A matter of principals}
Naturally, Team Alpha was quite upset about this. Ordonians took a stewardship approach out of concern that the event couldn't survive without two active teams. We gathered food, iron, blocks, and saplings from our collection to help them, then sailed to a small dock they called the embassy. It didn't take them long to recover. Their team quickly grew to a size of around six to eight active members compared to our core four, one of which was unable to participate for the second week. Some of us questioned if giving them the care package was a wise idea, stating that it removed from them a lot of necessary work and travel to get some of those supplies, and at our own expense. We were able to help them `cross things off a checklist'. Thus, we can see a conflict between two instincts and two principals emerge. Do we do what seems best for the event, or do we serve only the interests of our team? This is the first part of the Ordonian Calculus.\footnote{Social and systemic anticipation during The Systems Event by the individual members. Risk and failure model assessments unique to members of Ordonia.}

This point is also echoed in Harlow's `The Consequences of Poor Design', where it was noted that the typical Minecraft player, one that is not from the IMS, would be more concerned with getting ahead rather than the “idealistic purism” of “a niche online community, [like the IMS]”.~\citep{Harlow(2026)} It can also be seen how the divergence in development over time alongside homophily within the IMS led to diverging subtextual readings about the bounds and priorities of the event.

Ordonia has had a longer, more thorough period of exploration, even while both teams had the same time window. Ordonia grew during the exploration phase, causing Pheelix to take leadership much more seriously to handle the increasing player base. One that did not materialize as expected during the exploitation phase. We also had members that were active in the IMS broadly, one of which would eventually take the reins of the systems experiment from Turtle later on.

Alpha benefited from a rush of recruitment and a relaxation of requirements to participate in the event after a serious lag in recruitment and activity for most of the exploration phase. They had less formal ground to work with and some of their active members applied a different subtext to the event, one that would otherwise serve them in a typical factions-based SMP.

To lay out the principals, the primary principal is the IMS which includes the organizers of TSE. The other Principal is whichever team the player is on. These principals would presumptively work against each other because of what was left to subtext. Ordonia's core four were in a high demand, low information environment with its members sharing more responsibility and losing interest.

Alpha didn't carry the weight of planning or fulfilling large systems in quite the same sense. The way that Alpha formed by the opening of the server led to a team that was primarily focusing on the most legible aspects of the event, with the formal focuses typically left to whoever was the leader at the time.

The differences in styles of leadership would prove consequential in the way members – stakeholders in their respective teams – related to power and team organization. Power and Legitimacy was established with roles and elections during the exploration phase in Ordonia. Alpha held an election quite some time after the opening of the server. Alpha did not perceive a need for formal elected roles beyond the sole leadership.

\subsection{Diplomacy}
Scorvon is a player whose history and perspectives were sadly missing from Revs `Who Decides?' paper. His role as the leader of Team Alpha was as iconic as the power struggle between him and Phantom\_Boi, who eventually beat him in Alpha's election. He was a `strongman' of sorts, and an extremely effective leader who saw Alpha's swift early progression.

A member of Ordonia intercepted a written book on the mushroom island pertaining to Scorvon's leadership and the talks of war. The journal states that war was likely to take place for `entertainment value'. They were aware of the mismatched team size, and the presence of fully geared PvPers on their team. They wanted to declare when all Ordonians were active. This would conflict with statements from both Scorvon and Phantom\_boi that the majority of team Alpha did not want war. The journal confirms this later by stating that many in team Alpha were skeptical about Scorvon, leading to Phantom\_Bois' victory. In Revs' interviews, Rev was led to believe that there were no power struggles in Alpha, and Scorvon was not mentioned in the manuscript as of this moment.~\citep{Rev(2026)}

We became more deliberate in our interactions with Team Alpha, and observed patterns in the way they organized, kept the condition of their town, and tended to coordinate in small sub-groups or as individuals. There were complaints from Alphas that critical pieces of information, such as the state of treaties, weren't communicated well or at all. There was confusion about whether the majority of Team Alpha wanted war or wanted to avoid war. The most active members were looking for war and often conducted espionage even after we've come to agreements about what we expect before anyone comes to the opposite team's land. The easiest `catch' was when they required assistance getting out of our blocked Nether portals, a measure we took after one of us was killed in the Nether.

However, it wasn't always quite hostile, thanks to Hilo's efforts with the Tea Parties and her work on the mushroom island. There were also trade arrangements as Alpha was in need of a lot of building materials, and they had the only Pale Garden biome on the map. In order to get the few resin clumps they had, we needed to be transparent in why we wanted resin. The biggest point of the resin trade for Ordonia was to gather information. We were `anti-monopoly', and this incident would give us clues into how Alpha operates. Resin and its related products had its uses, but was never critically important. Nonetheless it was a monopoly and a somewhat strategic one. The early-game question was what their behavior would be towards more strategic monopolies – like the nether expansion granting either team access to the Netherite Upgrade Template.

Despite the trades and agreements, trust was still uneasy. While individuals within team Alpha had their own reputation, we were still cautious dealing with any member of Team Alpha. We would make sure that Ordonia wasn't deprived of the template by creating a rail line hidden low in the nether, covered to be less conspicuous where it crossed rivers and lakes of lava. We suspected a couple of their members killed an Ordonian they assumed had the template, but it was in an ender-chest.

\section{Agentic gameplay}
Perhaps the instinct to dominate or grief isn't incidental to this setting, but rather is invited by it. Minecraft is primarily built as a singleplayer game. While it can be played multiplayer, it does not have its own mechanisms built into itself that meaningfully alter the gameplay or add reasons or mechanisms to compete or cooperate with others. The systems that vanilla Minecraft provides are maximally legible and exploitable. It is easy for any experienced player to reach the end game. Minecraft isn't alone in this though. For example, Paradox grand strategy games in singleplayer modes are often built to reward accumulation and expansion, to `paint the map' is rewarded until your game becomes too boring or tedious to manage. Multiplayer lobbies often ban a list of metas or strategies. Expansion becomes negotiated among neighbors, and exploitation of your resources becomes the meta for multiplayer lobbies. Players move from a utilitarian and agentic style of play to one more oriented towards stewardship both internally and externally. This environment is still competitive, but it remains one bounded by rules and negotiations to the benefit of everyone in the lobby.

These kinds of multiplayer strategy sessions also invite not just a different set of extrinsic goals, but invite intrinsic ones as well. You're responding to a different environment that doesn't exist in singleplayer, your choices are going to be different, more intrinsic to what is happening between players than the meta of the game itself.

Minecraft is built differently because it does not structurally create channels of strategic behavior, yet all the same it is a singleplayer game, with a simple and appealing progression chain. Build a base, get diamonds, maybe villagers, go to the Nether, then beat the dragon. It is a binary checklist of things you either did or did not do, and to some that is part of the point of the game. The game doesn't prioritize taking a creative route, navigating roadblocks, sticking by your own principles, or anything beyond the binary checklist in any structural manner. It is a mostly blank slate devoid of expectations, and most players will be accustomed to this. This is Agentic Gameplay.

\subsection{The exceptions}
Hermitcraft is one example of a model of stewardship, even when storylines reflect agency or rivalries for part or whole of a series. Pranksters and Institutions like the Permit Office offer some antagonism to varying degrees; however, it is all within the bounded context that Hermitcraft provides. No players are put into silos for extended periods of time. Contradicting goals serve a creative and narrative-driven goal. One could imagine though that private back channels exist, where surprises are managed and boundaries are discussed. Agency itself is subsumed through stewardship.

In a Systems Event context, this could be intentionally replicated by `Red Teaming'.\footnote{a term used in cybersecurity, where ethical hackers on the behalf of an organization act as antagonists, using the tools and exploits of real antagonists} A system would be designed to be tested by another group that assumes an antagonistic role, then report their assessment. In cybersecurity, ethical hackers have the same skills and tools as malicious ones, but their goals are non-destructive. They operate with permission from the organization while following a code of conduct.

Hermitcraft and similar SMPs are quite systems oriented due to their long running seasons. Hermitcraft is composed of over twenty members, whose presence helps to multiply the efforts of the individual. There is no siloing of players.

Survival Multiplayer Servers that focus on PvP, particularly siloed events like the civilization events hosted by Ish, rely on balanced teams and legible, predetermined structured systems from the organizers. The Systems Event lacked the structural support for both systems and for PvP. Members of the TSE were left to their own subtextual understanding and failure modelling. For Ordonia, we largely lost sight of the goal as we were carrying a debt from our exploration phase that we had difficulty fulfilling, and we were managing more internal and diplomatic tasks.

\subsection{The question of war}
The question of war was something that wasn't fully clear for us. There wasn't much to go off of in official documentation of what war means, the stakes of war, the period, and should it be a bilateral or unilateral `declaration'. Was war just another testing ground, or was it part of the aim of the event in itself, as opposing factions are prone to do?

Applications were supposedly screened against people who were `too PvP focused', yet war wasn't taken off the table. It was another thing left for the `vanilla' or `archaic' feel – undefined and underbaked axioms that were counterproductive to the event. It was something you could do, if you want.

But the event couldn't support a war-based factional narrative. There was no hardcore mode, unbalanced starting conditions, and no guidelines for it. Early on, I suggested a `gentleman's war' at the last week of the event. That plan didn't get very far, as activity dwindled after the second week. Admittedly, the idea of a unilaterally declared war was something that seemed out of place in this event to most Ordonians, largely because of the imbalance in teams and the subtext we have interpreted for the event.

We Ordonians were aware of, or at least told, that a sizable portion of Team Alpha was interested in war `for entertainment value', and that some saw it as a valid zero-sum path towards victory by the metrics of The Systems Experiment. Individually, we went back and forth on how much we should interact with Alpha. Some members would be suspicious and then cool off, only for others to be similarly suspicious. While we felt that the people of team Alpha were fairly transparent actors, we couldn't always take what they said at face value.

\subsection{The temperature of the tea}
Hilo's role in the event stood out as the host of tea parties early on, where potions were labelled as `Osmanthus Tea'. She became one of the primary builders on the neutral mushroom island alongside Blaze. She was seen as a neutral party, yet still at least nominally a part of Alpha. We were aware that she was collecting information during the event, and that she had some connection with the Publishing Department.

We recognized that her positions towards both Principals were unique. When there was little reason for Ordonia to interact with Alpha, her tea parties were a diplomatic bridge that didn't need a formal excuse outside itself. It was an invitation to come to them, see what they have built and developed under diplomatic auspices.

Visitations slowly dropped off as a Visa system implemented by Scorvon came into effect. This was an early power-grab of his before the election. My first visa was signed by Phantom\_Boi during their power struggle, only for Scorvon to declare it invalid and request that I fill out another visa. It didn't take long for most cross-team activity to switch to the mushroom island, bypassing the need to adhere to antisocial visa policies, the principle of which were sometimes violated in either direction anyway.

`Neutral' positioning aside, we all answered to a team, but we understood Hilo answered to the IMS and the publishing department in a way we weren't quite sure of yet. She adopted a stewardly position, but we wouldn't be sure of what that could have implied should she or the organizers wanted to take the event in a specific direction. It's understood that a need for cross-team proximity was a given, and that we didn't quite understand how the organizers planned things, or if they have had much planning at all. We largely predicted the potential of Team Alpha as Agents rather than stewards or co-managers.

If we were to view things from the agentic perspective, with war being `entertainment', we couldn't realistically provide a satisfactory argument against that. If war was a binary checklist item without guidelines or expectations, the results would be predictably lopsided. So, it wasn't really about preventing war on its own merits, but making war potentially expensive. Hilo's early diplomatic approach made it easier for us to make connections, trade deals, and social reconnaissance on Team Alpha.

The diplomatic game became one of learning about the members of the other team, getting candid and one-on-one disclosures, and making trade deals. The problem with trade deals in this environment though is leverage; each team would desire to have some kind of leverage over the other, but leverage spent is leverage lost. The economy of information proved more valuable than most commodities in this kind of agency-driven environment, where guidelines and resources were thin.

\subsection{The fall-off}
War wasn't prevented by diplomatic activity alone, but largely through inactivity. Ordonia's systems proved difficult to implement with the scale and resources we ended up with, our city planning ambitions were poorly planned and mismatched. Some Ordonians deliberately took a break from the server, while some simply found other things more interesting, and the ones that stood by burned out.

Alpha, from what we could see, accomplished what they felt they could. The last two weeks of the event were mostly quiet from their side, with exception of the Nether expansion and the opening of the End. There was never a moment when enough Ordonians would be on that declaring a war could be worthwhile.

\section{Conclusion}
We didn't manage to produce lasting systems that could be shown off, there was nothing guiding us towards novel cooperative approaches with so few members. There was little for either team to take back to the IMS, because in this type of loose, vanilla environment, the most interesting thing is what the other team is doing. Where The Systems Event may have been productive was as a study of inter- and intra-team dynamics, and as a lesson of how not to create events.

What could be done for future events depends on the direction that the IMS intends to take them. Turtle seems to have drawn inspiration from different sources while trying to design for the broadest appeal, but without more deliberate design, clear communication and intentions, and a proactive consideration for players' time, we may see a repeat of this season's turn out.

It might be more prudent to risk over-design of future events, so that each season carries with it an identity, a set of unique problems or challenges, or conditions imposed, such that the focus becomes player versus environment, or player within the temporary society. Vanilla doesn't always mean perfectly vanilla or without guidelines, and anarchy vanilla has a particularly short set of possible outcomes.

Through the lens of Agency and Stewardship, we may begin to understand how player behavior relates to the environment they find themselves in. Even when interests don't directly align, all parties should be aware of the interests of each other. The Systems Event was set up with Three Principal Entities all acting within their own interests while performing a mental calculus of each party's intentions and efficacy.

It is up to organizers and server administrators to understand and anticipate the resources at their disposal. The system drafts turned into mere homework. The asynchrony of development and implementation cooled motivation and led to systems misaligned with the actual conditions of exploitation. Failure to set and clearly communicate guidelines, goals, and expectations created an environment of uncertainty that was counterproductive for the assumed goals of the event.

The design of the event did not clearly define its principles. We can somewhat surmise the intentions through the resources and statements that do exist; however, no true mission statement or theory of design has been made available outside of a points system. The structure of the event worked against both the resources it had available and its own design goals.

In conclusion, there should be a theory of Stewardship for the IMS. Any events, experiments, SMPs, or otherwise should have a clear scope and structure, with the understanding that emergent systems will develop downstream from structural systems. Minecraft itself is structural, `vanilla', `semi-vanilla', `anarchy', etc., are not neutral choices, nor are all or most servers in a category going to be the same. These choices are structural and worthy of scrutiny.

Minecraft Server Theory is still young, but has a solid framework thanks to creators like Matcha\_Tea, Turtlestate, KingGame, Pinecone, and others. We have only started on what will be a rich tradition for Minecraft as a platform. Player and staff/owner behavior and motivations are as, if not more, important than the systems themselves.

\section{Acknowledgements}
I would like to acknowledge Harlow for her initial feedback on the previous draft. This essay was difficult to write, but I hope it provides needed insight from a participant in the event. I would also like to acknowledge Rev's work and interviews as it provides an `overhead' view of the event that is professional and fair to both teams. This essay might not have been finished without either of them.

I had anticipated writing this essay since the end of the second week, yet it has had several redrafts and reorganizations. It is currently two weeks after the closing of the event, but after reorganizations and assessments of bias I hope that this article is worthy of the IMS.

Lastly, I want to acknowledge Turtle and Harlow again for their role in organizing the event.

\section{Declarations}
The author has declared a conflict of interest due to being on the Ordonian team. No conflict of interest has been declared by the reviewer and editor. No AI use has been declared.

    {
        \raggedright
        \begin{thebibliography}{99}
            \bibitem[Anderson(2024)]{Anderson(2024)}
            Anderson, E., Holdsworth, J., Kosinski, M. (2024). \textit{What is red teaming?} IBM. https://www.ibm.com/think/topics/red-teaming
            \bibitem[blazeflow(2026)]{blazeflow(2026)}
            blaze\_flow. 9 July, 2026. “IMS: The Systems Event” Discord Server. \textit{And your victory will be empty if we never finish our systems in the first place.} https://discord.com/channels/1496328885514539210/1521941424768024798/1524800867562295306
            \bibitem[Davis(1997)]{Davis(1997)}
            Davis, J. H., Schoorman, F. D., Donaldson, L. (1997). \textit{Toward a stewardship theory of management.} Academy of Management Review, 22(1), 20–47.
            \bibitem[Harlow(2026)]{Harlow(2026)}
            Harlow. 2026. “The Consequences of Poor Design: A Critique of the Systems Experiment.” Institute of Minecraft Studies, Substack. https://imsresearch.substack.com/p/the-consequences-of-poor-design-a
            \bibitem[Johnson(2011)]{Johnson(2011)}
            Johnson, S. (2011). GD column 17: \textit{Water finds a crack. Designer Notes.} https://www.designer-notes.com/game-developer-column-17-water-finds-a-crack/
            \bibitem[Rev(2026)]{Rev(2026)}
            Rev. 2026. “Who decides? Leadership, legitimacy, and load in competing teams: A revised analysis of the IMS Systems Event.” IMS Draft Manuscript.
            \bibitem[scorevon(2026a)]{scorevon(2026a)}
            scorevon. 9 July, 2026a. “IMS: The Systems Event” Discord Server. \textit{Ordaina had a whole city built banks schools etc, one week later they were bombed and everything was destoryed.} Retrieved 30 August, 2026, from https://discord.com/channels/1496328885514539210/1521941424768024798/1524800586317303870
            \bibitem[scorevon(2026b)]{scorevon(2026b)}
            scorevon. 9 July, 2026b. “IMS: The Systems Event” Discord Server. \textit{you cant win if ur systems are gone.} Retrieved 30 August, 2026, from https://discord.com/channels/1496328885514539210/1521941424768024798/1524800630743630035
            \bibitem[scorevon(2026c)]{scorevon(2026c)}
            scorevon. 9 July, 2026c. “IMS: The Systems Event” Discord Server. \textit{wins a win.} Retrieved 30 August, 2026, from https://discord.com/channels/1496328885514539210/1521941424768024798/1524800918091075696
            \bibitem[seetite(2026)]{seetite(2026)}
            seetite (1509229671235977338). 21 July, 2026. “Institute of Minecraft Studies” Discord Server. \textit{my thoughts: 1) I don't see a major issue with this. how are you going to do an event without the rp.} Retrieved 30 August, 2026, from https://discord.com/channels/1461415682586574850/1529095357894496527/1529110575412547855
            \bibitem[Turtle(2026)]{Turtle(2026)}
            Turtle. 21 July, 2026. “Critiques of Civ Experiments.” PDF. Retrieved 31 August, 2026, from https://discord.com/channels/1461415682586574850/1461415683882750007/1529095357894496527
            \bibitem[Turtle(nd)]{Turtle(nd)}
            Turtle. n.d. “TSE Video.” \textit{Google Docs}. Retrieved 31 August, 2026, from https://discord.com/channels/1461415682586574850/1461415683882750007/1529095357894496527
            \bibitem[Turtle(2026a)]{Turtle(2026a)}
            turtle2770. 7 June, 2026a. “IMS: The Systems Event” Discord Server. \textit{After careful consideration, I believe that there are actually two issues here even though both are not directly your question.} Retrieved 30 August, 2026, from https://discord.com/channels/1496328885514539210/1512536060142424084/1513132076097081505
            \bibitem[Turtle(2026b)]{Turtle(2026b)}
            turtle2770. 11 June, 2026b. “IMS: The Systems Event” Discord Server. \textit{Well I am gonna be a bit blunt tbh with some of this like Pearls and such. Since we are...} Retrieved 30 August, 2026, from https://discord.com/channels/1496328885514539210/1497330526238019584/1514767982998323333
            \bibitem[Turtle(2026c)]{Turtle(2026c)}
            turtle2770. 14 June, 2026c. “Institute of Minecraft Studies” Discord Server. \textit{The Systems Event is coming soon @everyone (I normally wouldn't, but the event members asked) In case any of you...} Retrieved 30 August, 2026, from https://discord.com/channels/1461415682586574850/1477119527203377182/1515884966477697094
        \end{thebibliography}
    }